\documentclass[12pt,a4paper]{article}
\usepackage[margin=1in]{geometry}
\usepackage{hyperref}
\usepackage{listings}
\usepackage{enumitem}
\usepackage{xcolor}

\lstset{
  basicstyle=\ttfamily\small,
  breaklines=true,
  frame=single,
  backgroundcolor=\color{gray!10}
}

\title{Prompt Log -- 2026-03-25}
\author{Swaroop Formulation Industries}
\date{2026-03-25}

\begin{document}
\maketitle

\section{Prompt Log Entry -- 001}

\begin{description}
  \item[Timestamp] 2026-03-25T20:30:00
  \item[Level] Overview
  \item[Model] claude-4.6-opus (Cursor AI)
  \item[Category] system-prompt
  \item[Prompt] \begin{lstlisting}
Follow Project_Report___Biodegradable_Packaging_Material.pdf and
_0_Resources/_certificates/. Suggest a bottom up and top down approach
for website development from beginner to expert level (scale 1 to 6).
Include chatbot, RAG, Agentic RAG, Adaptive RAG integration. Add Data
Analysis, dashboarding, certifications pages. Ideas for data scraping,
Airflow ETL, GenAI automation. Version control, cloud deployment hints.
Documentation in .md and .tex formats.
  \end{lstlisting}
  \item[Context] Initial project planning session to create a 6-level progressive website development roadmap for Swaroop Formulation Industries biodegradable packaging business.
  \item[Outcome] Generated comprehensive 6-level roadmap plan covering: L1 Foundation Chatbot, L2 RAG-Enabled Chatbot, L3 Agentic RAG + Certifications, L4 Adaptive RAG + Dashboards, L5 Full Platform + GenAI, L6 Enterprise K8s. Plan accepted and implementation started.
  \item[Files Modified] Plan file created, .cursor/rules/ (4 files), documentation/ (20+ files)
  \item[Tokens Used] \textasciitilde{}15000 input / \textasciitilde{}8000 output (estimated)
\end{description}

\end{document}
